\section{Cross-segment recall on real language}
\label{sec:language}

\subsection{Task}

A batch is a sequence of segments. The first segment states $n$ facts---key/value
pairs---embedded once; intervening segments are unrelated PG-19 prose that acts as
distractor filler; the final segment is a query that contains \emph{no} fact, only
a marker and one key, and asks the model to produce the corresponding value. Since
segments do not cross-attend and the query segment holds no facts, a correct answer
can only come from the store. We report exact-match recall and always run two
matched controls: store-disabled (the read is off) and scrambled-store (each
example reads a neighbour's store). Both must sit at chance for the number to mean
retrieval, and they do.

\subsection{From synthetic symbols to words}

On a synthetic multi-query associative recall task with single-symbol keys the
block reaches \recall{100} at 8 facts on three seeds, with both controls at
chance and the store reloaded in a fresh process at no loss. Moving to real
language exposed two obstacles.

\paragraph{The multi-token-key tax.} A fact key written in character-level text is
\emph{several} tokens, and the only token available at the retrieval point is the
shared delimiter; the read then becomes a uniform draw over the fact set, so
per-fact accuracy tracks $1/n$. Character-level recall consequently declines
monotonically with fact count (4 facts \recall{90}, 8 \recall{79}, 12 \recall{62},
13 \recall{55}) and caps around 13 facts. A single-token \emph{sub-word} (BPE)
key removes this tax: the key is one token again, competing directly with the
prose vocabulary. Cold, an 8-fact BPE model already recalls \recall{49.8} versus
\recall{10.3} for character keys---a $5\times$ gain with no per-character penalty.

\paragraph{The fact-count cliff is bootstrapping, not capacity.} Even with clean
keys, cold training exhibits a sharp cliff: recall is high at a few facts and
collapses to chance beyond a threshold, while the retrieval always lands on the
\emph{correct} slot (matching is intact; only the content is missing). Enlarging
the store does not help---the same recall is byte-identical at $2\times$ and
$4\times$ capacity---so the cliff is a gate/decoder bootstrapping deadlock, not a
capacity limit. A \emph{warm-start curriculum} escapes it: training a harder
fact-count from the checkpoint of an easier one ladders past the cliff. The
curriculum is the escape route for this single-site placement, but not the only
one---multi-layer placement clears the same cliff cold given a larger training
budget (\S\ref{sec:placement}), so the curriculum is best read as a
training-efficiency shortcut rather than a strict requirement.

\subsection{Results}

With single-token keys and the curriculum, the store recalls real-language facts
stated once and queried entirely out of context (\tabref{tab:language}):
\recall{100} at 8, 16 and 32 facts and \recall{99.1} at 64, replicated across two
seeds, with the matched scrambled control at chance ($0.16\%$) even at 64 facts.
Persistence holds on the language store: written, killed, reloaded in a fresh
process, recall equals in-process (\recall{100}); a mismatched store file falls to
$0\%$.

\begin{table}[t]
\centering
\begin{tabular}{lccccc}
\toprule
& 8 facts & 16 & 32 & 64 & scrambled (64) \\
\midrule
BPE + curriculum & \recall{100} & \recall{100} & \recall{100} & \recall{99.1} & $0.16\%$ \\
\bottomrule
\end{tabular}
\caption{Real-language recall: facts stated once in PG-19 prose, queried out of
context, answered from the store. Two seeds; matched scrambled control at chance.}
\label{tab:language}
\end{table}
